\documentclass[border=8pt]{standalone}
\usepackage{ctex}
\usepackage{tikz}
\usetikzlibrary{arrows.meta, positioning, calc}
\begin{document}
\begin{tikzpicture}[
  font=\small,
  box/.style={draw, rounded corners=3pt, minimum height=8mm, align=center, fill=#1},
  sbox/.style={box=gray!10, draw=gray!60, minimum width=30mm},
  dbox/.style={box=yellow!25, draw=yellow!70!black, minimum width=30mm},
  mbox/.style={box=blue!12, draw=blue!50!black, minimum width=30mm},
  rbox/.style={box=green!12, draw=green!50!black, minimum width=30mm},
  cbox/.style={box=orange!15, draw=orange!70!black, minimum width=30mm},
  arr/.style={-{Stealth[length=2.2mm]}, thick},
]

% 问题四 女胎异常判定流程
\node[sbox] (inp) {输入：女胎数据\\（Z 值、GC 含量、读段数、BMI 等）};
\node[dbox, below=8mm of inp] (base) {基线方法\\$|Z_{13/18/21}|>3$ 阈值判定};
\node[mbox, below=8mm of base] (feat) {特征构造\\X 染色体浓度 + 染色体 Z 值 +\\GC 含量 + 读段比例 + BMI + 年龄};
\node[mbox, below=8mm of feat] (model) {分类模型\\逻辑回归 / 随机森林\\（5 折分层交叉验证）};
\node[cbox, below=8mm of model] (eval) {模型评估\\ROC / AUC / 混淆矩阵 / 灵敏度特异度};
\node[rbox, below=8mm of eval] (out) {输出：女胎异常判定方法\\+ 与基线对比};

\draw[arr] (inp) -- (base);
\draw[arr] (base) -- (feat);
\draw[arr] (feat) -- (model);
\draw[arr] (model) -- (eval);
\draw[arr] (eval) -- (out);

\end{tikzpicture}
\end{document}
